% ===== 原文 pp.4–6（PDF p.14–16） =====

% [原文 p.4]

配有三个独立的复结构 $I,J,K$ 与三个独立的辛形式 $\omega_I,\omega_J,\omega_K$。因此，
视我们选取 $I,J,K$ 中的哪一个，$\cM_H(\LG,C)$ 可以被看作复流形（相应地，辛流形）。
与 $t=1$ 约化相关的辛形式是 $\omega_K$。（关于 Hitchin 模空间的更多内容，见第 8 章。）

在 B-侧，$t=\ii$ 的扭转理论约化为以 $\cM_H(G,C)$ 为靶空间的拓扑 B-模型，此时它是把
$\cM_H(G,C)$ 看作复结构 $J$ 下的复流形。对偶性的这一二维表现形式——在一对靶空间上的
A-模型与 B-模型之间——在物理学家中被称为\emph{镜像对称} [22]。为了与几何朗兰兹建立
联系 [1]，人们取 $\Sigma$ 为带边界的黎曼面，并考虑沿边界的每个连通分支可以施加在
2d 理论各场上的所有可能的边界条件的集合。在 A-侧，物理学家把这些边界条件称为
\emph{A-膜}，其中包括装备了平坦线丛的靶空间的拉格朗日子流形。（二维理论的膜在第 9 章
有广泛讨论。）在 B-侧，边界条件被称为 \emph{B-膜}，其中包括靶空间的复子流形。

% [Tutor Note: p.4 “A-膜 = 拉格朗日子流形 + 平坦线丛”是 A-模型 BPS 边界条件在几何极限下的
%  经典描述（Kapustin–Orlov）；一般情形下 A-膜是（可能带平坦 $U(1)$ 丛的）拉格朗日子流形，
%  但在 $\cM_H$ 这种非紧超 Kähler流形上还需附加“余迷向”（coisotropic）条件的推广，
%  见 [56] 以及第 7 章对修正 A-模型边界条件的讨论。]

因此，S-对偶蕴含如下等价：
\begin{equation}
  \eqbox{A-膜 \\ 在 $\cM_H(\LG,C)$ 上，\\ 辛形式 $\omega_K$}
  \;\longleftrightarrow\;
  \eqbox{B-膜 \\ 在 $\cM_H(G,C)$ 上，\\ 复结构 $J$} \tag{1.2}
\end{equation}

(1.1) 的 B-侧上的平坦 $G_{\C}$-丛对应于复结构 $J$ 下 $\cM_H(G,C)$ 的点。这些点又标记了
一类特殊的 B-膜，即支撑在单个点上的 0-膜。因此，我们认出 (1.1) 的 B-侧上的对象是
(1.2) 的 B-侧上对象的一个子类。给定支撑在 $p\in\cM_H(G,C)$ 上的一个 0-膜，
Strominger、Yau 与 Zaslow 给出了一个方案 [19]，说明如何在镜像对称下求 $p$ 的像：
它被映到带有平坦酉线丛、绕在 $\cM_H(\LG,C)$ 的某个中维子流形上的一个 A-膜，
这个子流形就是 $p$ 的所谓\emph{对偶 Hitchin 纤维}。Kapustin 与 Witten 通过说明这个
A-膜如何被映到 $\Bun_{\LG}C$ 上的一个 Hecke 特征层，完成了论证。

% [Tutor Note: p.4 式 (1.2) 的核对：SYZ 镜像对称把 $(X,\omega)$ 上的 A-膜与镜像
%  $(\check X,J)$ 上的 B-膜对应起来；在 Hitchin 系统的情形，镜像变换是沿 Hitchin 纤维化
%  $\cM_H\to\cB$ 的纤维做 T-对偶，故 $\cM_H(G,C)$ 中一点 $p$（0-膜）的镜像是
%  $\cM_H(\LG,C)$ 中对偶 Hitchin 纤维（一个复环面 $T^{2n}$，$n=(g-1)\dim G$）上的
%  A-膜。该纤维恰是拉格朗日的（在 $\omega_K$ 下），这与 A-膜的条件一致。]

事实上，(1.2) 在两个意义上比 (1.1) 体现了更丰富的等价。第一，(1.2) 中的方框包含比
(1.1) 更多的对象：除了 0-膜之外还有多得多的 B-膜。但或许从比较数学结构的观点看更有
意思的是：(1.2) 的等价不只是集合之间的等价，而是\emph{范畴}之间的等价。范畴是一个
对象集合，连同态射的汇集：连接这些对象并满足复合性质的箭头。（膜范畴的态射对应于
坐在两个边界条件交接处的局部算子，我们在第 4 章讨论。）因此，(1.2) 的等价是一种
特殊的映射，它尊重范畴的态射所体现的“刚性”：即范畴之间的一个函子。

% [原文 p.5]

\section{Kapustin--Witten 工作的推广}

在本论文中，我们沿几个方向推广 [1]。把在黎曼面上的紧化这一概念加以推广，我们在第 5 章
改为考虑在\emph{圆}上紧化到三维有效 TFT。在 A-侧，我们将发现 $t=1$ 的理论紧化为一个
鲜有研究的 3d TFT，称为\emph{带规范场的 3d A-模型}。3d A-模型 TFT 是一个三维拓扑
$\sigma$-模型，可对任意黎曼靶流形 $X$ 定义；它首先在 [6] 中对一般靶空间作了研究
（此处 §B.2 有回顾），可以被看作 2d A-模型的三维类比。当靶空间容许某个李群的作用时，
它可以与一个 3d 拓扑规范理论耦合，构成带规范场的拓扑 $\sigma$-模型，我们将在第 3 章
更详细地讨论。（带规范场的拓扑 $\sigma$-模型将在本论文中居于核心地位，我们在附录 C
回顾其若干数学方面。）3d $\sigma$-模型的靶空间是群流形 $\LG$，其上配有它到自身的共轭作用。

在 B-侧，我们很幸运：3d 紧化后的理论更为熟悉，它是 Rozansky--Witten（RW）3d $\sigma$-模型的
带规范场版本（相关的规范场化写在 [4] 的附录中；我们在 §B.5 回顾）。靶空间是复化群 $G_{\C}$
的余切丛的全空间 $T^*G_{\C}$，其上配有 $G$ 在底空间 $G_{\C}$ 上的共轭作用以及在纤维上
诱导的余伴随作用。实际上，把纤维指定“鬼数”为二、把底空间指定为零会更方便；为表明这一点，
我们把靶空间记作 $T^*[2]G_{\C}$。

对偶性预言这两个 3d 带规范场的拓扑 $\sigma$-模型之间存在高度不平凡的等价
（而且，除了 [4] 与本论文的分析之外，完全未被触及）。特别地，它们的边界条件集合
应当同构：
\begin{equation}
  \eqbox{带规范场的 3d A-模型的\\边界条件，靶空间 $\LG$}
  \;\longleftrightarrow\;
  \eqbox{带规范场的 RW 模型的\\边界条件，靶空间 $T^*G_{\C}$} \tag{1.3}
\end{equation}

这一关于朗兰兹对应的“3d TFT”视角比 (1.2) 的 2d TFT 视角更为丰富。正如我们将在第 4 章
解释的，3d TFT 的边界条件具有 2-范畴的结构，包括对象、态射以及态射之间的态射
（即 2-态射）。概括地说：(1.1) 是对象集合的等价，而 (1.2) 是范畴（对象 + 态射）的等价，
(1.3) 则是 2-范畴（对象 + 态射 + 2-态射）的等价。在这一进程的每一步，人们都加入了
“额外的”对象，并把上一步中只是隐含的刚性显式化。事实上，(1.3) 当然还不是终点：
由于它们的四维起源，3d 边界条件“暗中”具有一个被对偶性保持的辫状、幺半结构；
不过本论文止步于 2-范畴的层次，把这一辫状、幺半结构留给未来的工作。

% [Tutor Note: p.5 关于“四维起源使 3d 边界条件带有辫状幺半结构”：直观地说，
%  4d 理论中的面算子可以沿两个横向方向彼此靠近、堆叠（融合），这给出幺半结构（张量积）；
%  而在四维中两个横向的二维曲面可以相互绕行（在 3d 约化中表现为边界条件的交换），
%  因而张量积一般不对易而只满足辫状关系。本论文不展开这一点，作者在 [1, 4] 中有相关讨论。
%  % 批注脚注：本批注不含数学依赖，仅提示“braided monoidal”一词的来源。]

% [原文 p.6]

我们推广 [1] 的另一个方向，是考虑比 [1] 中的线算子与 2d 膜更广的一类可观测量。
在第 6 章，我们对支撑在余维二子流形上的算子，即\emph{面算子}，进行分类。
应用于几何朗兰兹的最简单的一类面算子由 S. Gukov 与 E. Witten [2] 研究过；
不过我们将发现，扭转理论中还有更多得多的、更一般的面算子。对于 $t=\ii$ 的理论，
我们得到面算子 2-范畴的完整描述；对于其它扭转理论，我们只对阿贝尔规范群得到描述——
这正是（上面已指出的）那个事实的一个例证：B-侧的分析总比 A-侧更受控制。

事实上，向面算子的推广与向圆紧化的推广密切相关：因为借助第 4 章讨论的维数约化技巧，
4d TFT 的面算子 2-范畴恰好等同于其 3d 约化的边界条件 2-范畴。

% [Tutor Note: p.6 “维数约化技巧”的精确含义：把 4d 理论放在 $T^2\times$（其余）上沿其中一个
%  圆周紧化，靠近某条边界的 4d 面算子会约化为 3d 理论中的边界条件；反过来，给定 3d 理论的一个
%  边界条件，沿第四维方向平移即得到 4d 理论中沿该方向延伸的面算子。这一对应在第 4 章 §4.4 与
%  第 6 章中反复使用。]

此外，我们还会把另外两个可由 $\cN=4$ 经扭转程序得到的 4d TFT 纳入这个故事：$t=0$ 的
扭转理论，以及另一个（不属于同一族的）理论——由于它在 [3] 中所起的作用，我们称之为
\emph{Vafa-Witten（VW）扭转理论}。这些理论是自对偶的（只要适当地倒转耦合并把规范群
换为其朗兰兹对偶群）。$t=0$ 理论的某些方面也在 [35] 中被考虑过，它本质上等价于带非零
瞬子项的 $t=1$ 理论（该情形也在 [1] 中被考虑过）。已经证明，它为所谓的\emph{量子几何
朗兰兹对应}（(1.1) 的一个涉及量子参数的推广）提供了自然的框架。

在第 3 章，我们将定义一个全新的 TFT，即带规范场的 4d A-模型（与带规范场的 3d A-模型
密切类似），并说明 Vafa-Witten 理论可以看作靶空间取为规范群的李代数 $\fg$ 的带规范场的
4d A-模型。在第 5 章我们将说明，$t=0$ 理论与 Vafa-Witten 理论在 $S^1$ 上紧化为同一个
3d TFT：以 $\fg$ 为靶空间的带规范场的 3d A-模型的某个版本。我们还将得到 $t=\ii$ 与
$t=1$ 理论在圆上紧化的三维描述。

在第 4 章，我们将更详细地解释为什么 2d TFT 的边界条件具有范畴的结构，而 3d TFT 的
边界条件具有 2-范畴的结构。在第 6 章，我们将分析 3d 理论的边界条件，等价地，4d 理论的
面算子。在第 7 章我们构造\emph{修正 A-模型}，这是一个可对带有相容 $U(1)$ 作用的一般
Kähler 靶空间定义的 2d TFT。在第 8 章，我们说明 Vafa-Witten 理论在黎曼面 $C$ 上紧化为
以 $\cM_H(G,C)$ 为靶空间的修正 A-模型。

最后，在第 9 章，我们将得到 2d 玻色 $\sigma$-模型以及具有不同超对称量的 2d $\sigma$-模型的
拓扑缺陷的几何分类。（这最后一个专题多少处于本论文主线之外。）
